% ============================================================
%  PLANTILLA ANTEPROYECTO – INVESTIGACIÓN APLICADA
%  Corporación Universitaria Comfacauca – Unicomfacauca
%  Normas APA 7.ª edición
%
%  CÓMO COMPILAR (orden recomendado en terminal):
%    1. pdflatex  anteproyecto.tex
%    2. biber     anteproyecto
%    3. pdflatex  anteproyecto.tex
%    4. pdflatex  anteproyecto.tex   (segunda pasada para referencias cruzadas)
%
%  Si tu editor es Overleaf, simplemente haz clic en "Recompile".
% ============================================================

% ------------------------------------------------------------
%  CLASE DEL DOCUMENTO
%  - 12pt  : tamaño de fuente base (APA recomienda 11-12 pt)
%  - letterpaper : tamaño carta
%  - oneside     : impresión de un solo lado
% ------------------------------------------------------------
\documentclass[12pt, letterpaper, oneside]{article}


% ============================================================
%  PAQUETES
% ============================================================

% --- Codificación e idioma ---
\usepackage[utf8]{inputenc}       % Permite escribir tildes directamente
\usepackage[T1]{fontenc}          % Fuentes con caracteres latinos
\usepackage[spanish]{babel}       % Idioma español (seccionado, fechas, etc.)

% --- Fuente: Times New Roman (requerida por APA) ---
\usepackage{mathptmx}             % Equivalente a Times New Roman en LaTeX

% --- Márgenes APA 7: 2.54 cm (1 pulgada) en todos los lados ---
\usepackage[
  top=2.54cm,
  bottom=2.54cm,
  left=2.54cm,
  right=2.54cm
]{geometry}

% --- Interlineado ---
\usepackage{setspace}
\doublespacing                    % Doble espacio (norma APA 7)

% --- Párrafos: sangría de primera línea ---
\usepackage{indentfirst}          % Sangría en primer párrafo de cada sección
\setlength{\parindent}{1.27cm}    % 0.5 pulgadas = 1.27 cm (APA 7)
\setlength{\parskip}{0pt}         % Sin espacio extra entre párrafos

% --- Encabezado y pie de página ---
\usepackage{fancyhdr}
\pagestyle{fancy}
\fancyhf{}                        % Limpia encabezado y pie predeterminados
\fancyhead[R]{\thepage}           % Número de página arriba a la derecha (APA 7)
\renewcommand{\headrulewidth}{0pt}% Sin línea decorativa en el encabezado

% --- Hipervínculos (URLs en referencias) ---
\usepackage[
  hidelinks,                      % Sin recuadros de color en los enlaces
  pdftitle={Título del Anteproyecto},
  pdfauthor={Tu Nombre}
]{hyperref}
\usepackage{url}

% --- Figuras y tablas ---
\usepackage{graphicx}             % Para insertar imágenes con \includegraphics
\usepackage{booktabs}             % Tablas con estilo profesional
\usepackage{caption}              % Personalizar leyendas de figuras/tablas
% Formato APA de leyenda: "Nota." en cursiva, tamaño 10 pt
\captionsetup{
  font=small,
  labelfont=it,
  labelsep=newline,
  justification=raggedright,
  singlelinecheck=false
}

% --- Listas ---
\usepackage{enumitem}             % Control fino de listas con \begin{itemize}

% --- Referencias con biblatex + biber (APA 7) ---
\usepackage[
  style=apa,          % Estilo APA automático
  backend=biber,      % Motor de bibliografía (más moderno que bibtex)
  sortcites=true,
  sorting=nyt         % Ordena por Nombre, Año, Título
]{biblatex}
\addbibresource{referencias.bib}  % Vincula tu archivo .bib


% ============================================================
%  INICIO DEL DOCUMENTO
% ============================================================
\begin{document}


% ------------------------------------------------------------
%  PORTADA
%  (La portada no lleva número de página)
% ------------------------------------------------------------
\thispagestyle{empty}             % Sin encabezado ni pie en la portada

\begin{center}
  \vspace*{2cm}

  {\large \textbf{Título del Anteproyecto}} \\[0.5cm]  % <-- Cambia este título

  \vspace{1.5cm}

  {\normalsize Autores: Rodrigo Andres Realpe, Jose Luis Valencia Avirama} \\[0.3cm]    % <-- Tu nombre
  {\normalsize Director: Nombre del Director} \\[0.3cm]
  {\normalsize Programa: Ingeniería de Sistemas} \\[0.3cm]
  {\normalsize Universidad: Corporación Universitaria Comfacauca} \\[0.3cm]
  {\normalsize Popayán, Colombia} \\[0.3cm]
  {\normalsize \the\year}                               % Año actual automático

\end{center}

\newpage


% ------------------------------------------------------------
%  TABLA DE CONTENIDO
%  LaTeX la genera automáticamente con \tableofcontents.
%  Se actualiza sola cada vez que compilas.
% ------------------------------------------------------------
\thispagestyle{empty}
\tableofcontents
\newpage

% Desde aquí empieza la numeración de páginas
\setcounter{page}{1}


% ------------------------------------------------------------
%  RESUMEN EJECUTIVO / ABSTRACT
% ------------------------------------------------------------
\begin{abstract}
  % Escribe aquí un párrafo de 150–250 palabras que resuma:
  % problema, objetivo, metodología y resultados esperados.
  Escribe aquí tu resumen ejecutivo. Debe tener entre 150 y 250 palabras
  y cubrir: el problema que aborda el proyecto, el objetivo general,
  la metodología empleada y los resultados que se esperan obtener.
  \\
  
  \begin{center}
    \textbf{Palabras clave: }\textit{Aquí van las palabras clave, separdas por comas}
  \end{center}
\end{abstract}

\newpage


% ============================================================
%  CAPÍTULO 1 – INTRODUCCIÓN Y PLANTEAMIENTO
% ============================================================
\section{Introducción y Planteamiento del Problema}
% \section{} crea un capítulo numerado (1, 2, 3…).
% \subsection{} crea una sección dentro del capítulo (1.1, 1.2…).
% \subsubsection{} crea una subsección (1.1.1, 1.1.2…).
% Todos aparecen automáticamente en la tabla de contenido.

\subsection{Introducción}
Escribe aquí el texto introductorio del proyecto.

\subsection{Pregunta de Investigación}
¿Cuál es la pregunta central que orienta esta investigación?

\subsection{Planteamiento del Problema}
Describe aquí el problema que motiva el proyecto. Puedes citar fuentes
con \parencite{apellido2024} (formato APA entre paréntesis) o con
\textcite{apellido2024} (formato APA integrado en la oración).

\subsection{Justificación}
Explica por qué es importante resolver este problema y qué valor aporta
la investigación para la institución, la comunidad o la industria.

\subsection{Objetivos}

\subsubsection{Objetivo General}
Redacta aquí el objetivo general del proyecto. Inicia con un verbo en
infinitivo (p.\,ej., Diseñar, Desarrollar, Implementar).

\subsubsection{Objetivos Específicos}
% Lista numerada con \begin{enumerate}
\begin{enumerate}[label=\arabic*.]   % Numeración 1. 2. 3.
  \item Primer objetivo específico.
  \item Segundo objetivo específico.
  \item Tercer objetivo específico.
  % Agrega o elimina ítems según necesites
\end{enumerate}

\subsection{Alcance y Limitaciones}
Define hasta dónde llega el proyecto y cuáles son las restricciones o
condiciones que acotan su desarrollo.

\subsection{Hipótesis}
% Esta sección es opcional. Si no aplica, puedes eliminarla o comentarla
% poniendo % al inicio de cada línea.
Plantea aquí la hipótesis del proyecto (cuando aplique).
\parencite{hijji_multivocal_2021}

\newpage


% ============================================================
%  CAPÍTULO 2 – MARCO TEÓRICO
% ============================================================
\section{Marco Teórico}

\subsection{Antecedentes de la Investigación}
Presenta investigaciones o proyectos previos relevantes para tu tema.

\subsection{Estado del Arte}
Describe el nivel de avance actual del conocimiento sobre el problema
que abordas \parencite{autor2022}.

\subsection{Fundamentación Teórica}
Define y desarrolla los conceptos clave que sustentan el proyecto.


\newpage


% ============================================================
%  CAPÍTULO 3 – ANÁLISIS DEL PROBLEMA
% ============================================================
\section{Análisis del Problema}

\subsection{Análisis del Problema}
Desarrolla aquí el análisis detallado del problema identificado.

\subsection{Árbol de Problemas}
% Para insertar una imagen usa:
%   \includegraphics[width=\textwidth]{nombre_del_archivo.png}
% Guarda la imagen en la misma carpeta que este archivo .tex.
% Ejemplo:
%
% \begin{figure}[h]
%   \centering
%   \includegraphics[width=0.85\textwidth]{arbol_problemas.png}
%   \caption{\textit{Nota.} Descripción del árbol de problemas.}
%   \label{fig:arbol_problemas}
% \end{figure}

[Inserta aquí el diagrama del árbol de problemas]

\subsection{Árbol de Objetivos}
[Inserta aquí el diagrama del árbol de objetivos]

\newpage


% ============================================================
%  CAPÍTULO 4 – METODOLOGÍA
% ============================================================
\section{Metodología}

\subsection{Tipo de Estudio / Diseño de Investigación}
Describe el enfoque metodológico (cuantitativo, cualitativo, mixto) y
el tipo de investigación (exploratoria, descriptiva, experimental, etc.).

\subsection{Población y Muestra}
Caracteriza la población de estudio y explica el proceso de muestreo.

\subsection{Técnicas e Instrumentos de Recolección de Datos}
Describe las técnicas utilizadas (entrevistas, encuestas, observación,
análisis documental, etc.) y los instrumentos asociados.

\subsection{Procedimiento}
Detalla paso a paso cómo se llevará a cabo la investigación.

\subsection{Análisis de Datos}
Explica cómo procesarás e interpretarás los datos recolectados.

\newpage


% ============================================================
%  CAPÍTULO 5 – PLAN DE TRABAJO / CRONOGRAMA
% ============================================================
\section{Plan de Trabajo y Cronograma}

Describe las fases y actividades del proyecto. Puedes insertar aquí una
tabla o diagrama de Gantt como figura (ver instrucciones en el
Capítulo~3 para insertar imágenes).

[Inserta aquí la tabla de actividades con fechas y responsables]

\newpage


% ============================================================
%  CAPÍTULO 6 – PRESUPUESTO
% ============================================================
\section{Presupuesto}

\subsection{Recursos Materiales}
Lista los materiales, equipos y licencias necesarios para el proyecto.

\subsection{Recursos Humanos}
Describe el equipo de trabajo y sus roles.

\subsection{Costos Estimados}
[Inserta aquí la tabla de costos estimados]

\newpage


% ============================================================
%  CAPÍTULO 7 – MATRIZ DE RIESGOS
% ============================================================
\section{Matriz de Riesgos}

Identifica los principales riesgos del proyecto, su probabilidad de
ocurrencia, su impacto potencial y las estrategias de mitigación.

[Inserta aquí la tabla de matriz de riesgos]

\newpage


% ============================================================
%  CAPÍTULO 8 – RESULTADOS ESPERADOS E IMPACTO
% ============================================================
\section{Resultados Esperados e Impacto}

\subsection{Productos o Entregables Esperados}
% Lista de viñetas con \begin{itemize}
\begin{itemize}
  \item Primer entregable esperado.
  \item Segundo entregable esperado.
  % Agrega o elimina ítems según necesites
\end{itemize}

\subsection{Beneficios para la Institución, Comunidad o Industria}
Describe el impacto positivo que tendrá el proyecto al concluirse.

\newpage


% ============================================================
%  CAPÍTULO 12 – CRONOGRAMA DE IMPLEMENTACIÓN
%  (Numeración según la guía institucional)
% ============================================================
\section{Cronograma de Implementación}

[Inserta aquí el diagrama de Gantt u otra representación gráfica]

\newpage


% ============================================================
%  REFERENCIAS BIBLIOGRÁFICAS
%  biblatex las genera automáticamente desde referencias.bib.
%  Solo aparecen las fuentes que hayas citado en el texto.
% ============================================================
\section{Referencias Bibliográficas}
\printbibliography[heading=none]  % "heading=none" evita el título doble

\newpage


% ============================================================
%  ANEXOS
%  Etiquétalos como Anexo A, Anexo B, etc. (norma APA 7).
%  Menciona cada anexo en el texto antes de presentarlo.
% ============================================================
\section*{Anexos}              % El * evita que "Anexos" aparezca numerado
\addcontentsline{toc}{section}{Anexos}  % Lo agrega a la tabla de contenido

\subsection*{Anexo A – Nombre del Primer Anexo}
Descripción o contenido del primer anexo.

\subsection*{Anexo B – Nombre del Segundo Anexo}
Descripción o contenido del segundo anexo.


% ============================================================
%  FIN DEL DOCUMENTO
% ============================================================
\end{document}